% content/01_introduction.tex
\newpage
\phantomsection%
\section*{Введение}
\addcontentsline{toc}{section}{Введение}

База данных общежития предназначена для учёта студентов, мест проживания, заселений и платежей. При регистрации заселения сведения о студенте связываются с выделенным ему местом, а каждый платёж относится к соответствующему заселению. Для подготовки сводок по студентам требуется объединять данные из нескольких связанных таблиц.

В работе на примере PostgreSQL рассматриваются представления, триггеры, роли пользователей, хранимые функции и процедуры. Эти средства позволяют формировать сводки, поддерживать актуальность расчётных данных и выполнять часть операций учёта на уровне СУБД. Отдельно рассматривается влияние изоляции транзакций на результаты чтения при одновременной работе пользователей.

В результате реализовано автоматическое обновление счётчиков заселений и платежей, разработана процедура регистрации заселения вместе с платежом и разграничены права на просмотр сводки и изменение исходных записей. В сценарии с двумя параллельными транзакциями продемонстрировано неповторяющееся чтение суммы платежа на уровне изоляции \texttt{READ COMMITTED}.
